\subsection{The affine nonsqueezing theorem}

Define an {\emphcolor affine symplectomorphism} of $\bR^{2n}$ to be a map $\psi\colon\bR^{2n}\to\bR^{2n}$ of the form
$$ \psi(z) = \Psi(z) + z_0,\qquad \Psi\in\Sp(2n),\ z_0\in\bR^{2n} . $$
Let $\ASp(2n)$ denote the set of affine symplectomorphisms, this is a group as
$$ \psi\circ\psi'(z) = \Psi\circ\Psi'(z) + \Psi(z'_0) + z_0,\qquad \psi^{-1}(z) = \Psi^{-1}(z) - \Psi^{-1}(z_0) $$
are affine symplectomorphisms.

The affine nonsqueezing theorem says that a ball in $\bR^{2n}$ can only be embedded into a symplectic cylinder if the ball's radius is smaller than the cylinder.
If we let $e_1,\dots,e_n,f_1,\dots,f_n$ denote the standard basis of $\bR^{2n}$, then the symplectic cylinder of radius $R>0$ is
$$ Z^{2n}(R) = B^2(R) \times \bR^{2n-2} = \set{z\in\bR^{2n}}[\iprod{z,e_1}^2+\iprod{z,f_1}^2\leq R^2] \subseteq \bR^{2n} . $$

\bthrm[title={The affine nonsqueezing theorem}, name={affnonsqueeze}]

    Let $\psi\in\ASp(2n)$ be an affine symplectomorphism such that $\psi(B^{2n}(r))\subseteq Z^{2n}(R)$.
    Then $r\leq R$.

\ethrm

\bproof

    We may assume that $r=1$ and let $\psi(z)=\Psi(z)+z_0$.
    Then define
    $$ u = \Psi^\top e_1,\qquad v = \Psi^\top f_1,\qquad a = \iprod{e_1,z_0},\qquad b = \iprod{f_1,z_0} . $$
    Because $\Psi\in\Sp(2n)$ we have
    $$ \omega_0(u,v) = \omega_0(\Psi^\top e_1,\Psi^\top f_1) = \omega_0(e_1,f_1) = 1 . $$
    The condition $\psi(B^{2n}(1))\subseteq Z^{2n}(R)$ is equivalent to
    $$ \sup_{\norm z=1}\bigl(\iprod{e_1,\psi(z)}^2 + \iprod{f_1,\psi(z)}^2\bigr) = \sup_{\norm z=1}\bigl((\iprod{u,z}+a)^2 + (\iprod{v,z}+b)^2\bigr) \leq R^2 . $$
    Because $1=\omega_0(u,v)\leq\norm u\norm v$ (as $\omega_0(x,y)=\iprod{x,J_0y}$ and $J_0$ is norm-preserving), one of $u,v$ must have norm at least $1$.
    Without loss of generality $\norm u\geq1$, and then let $z=\pm u/\norm u$ whose sign agrees with $a$'s.
    Then
    $$ (\iprod{u,z} + a)^2 \geq \norm u^2 \geq 1 , $$
    which then implies that $1\leq R^2$ and so $1\leq R$ as desired.
    \qed

\eproof

Because $B^{2n}(R)\subseteq Z^{2n}(R)$, we have that
$$ R = \sup\set{r>0}[\hbox{there exists an affine/linear symplectomorphism $\psi$ such that $\psi(B^{2n}(r))\subseteq Z^{2n}(R)$}] . $$
Either affine or linear in the above supremum gives the same result.

Call a set $B\subseteq\bR^{2n}$ an {\emphcolor linear} (resp.\ {\emphcolor affine}) {\emphcolor symplectic ball} of radius $r$ if it is linearly (resp.\ affinely) symplectomorphic
to $B^{2n}(r)$.
Similarly $Z\subseteq\bR^{2n}$ is a {\emphcolor linear} (resp.\ {\emphcolor affine}) {\emphcolor symplectic cylinder} of radius $R$ if it is linearly (resp.\ affinely symplectomorphic
to $Z^{2n}(R)$.

The radius of a symplectic ball or cylinder is invariant.
This follows from the above theorem: if $B$ is symplectomorphic to $B^{2n}(r)$ and $B^{2n}(r')$, then $B^{2n}(r),B^{2n}(r')$ are symplectomorphic.
By the above theorem this means that $r=r'$.
And the radius of a symplectic cylinder is
$$ R = \sup\set{r>0}[\hbox{there exists an affine/linear symplectomorphism $\psi$ such that $\psi(B^{2n}(r))\subseteq Z^{2n}(R)$}] . $$

A matrix $\Psi\in\Mat_{2n}$ has the {\emphcolor linear nonsqueezing property} if for every linear symplectic ball $B$ of radius $r$ and linear symplectic cylinder $Z$ of radius $R$
$$ \Psi B\subseteq Z \implies r\leq R . $$
By \refmath{affnonsqueeze} any symplectic matrix has the linear nonsqueezing property.
Call a matrix {\emphcolor anti-symplectic} if $\Psi^*\omega_0=-\omega_0$.
\refmath{affnonsqueeze} clearly holds for anti-symplectic matrices as well.

\bthrm[title={The affine rigidity theorem}, name={affrigid}]

    Let $\Psi\in\Mat_{2n}$ be nonsingular such that $\Psi,\Psi^{-1}$ have the linear nonsqueezing property.
    Then $\Psi$ is either symplectic or anti-symplectic.

\ethrm

\bproof

    Suppose $\Psi$ is neither symplectic nor anti-symplectic.
    Then neither is $\Psi^\top$ so there are $u,v\in\bR^{2n}$ such that
    $$ \omega_0(\Psi^\top u,\Psi^\top v) \neq \pm\omega_0(u,v) . $$
    We may assume that neither side is zero, and potentially by replacing $\Psi$ with $\Psi^{-1}$, $\abs{\omega_0(\Psi^\top u,\Psi^\top v)}<\abs{\omega_0(u,v)}$.
    We can rescale $u$ so that $\omega_0(u,v)=1$ so that
    $$ 0 < \lambda^2 = \abs{\omega_0(\Psi^\top u,\Psi^\top v)} < \omega_0(u,v) = 1 . $$
    Thus $\lspan\set{u,v}$ is a symplectic subspace of which $u,v$ is a symplectic basis, so we can extend it to a symplectic basis of all of $\bR^{2n}$ by considering
    its symplectic complement.
    So let $u_1,v_1,\dots,u_n,v_n$ be a symplectic basis of $\bR^{2n}$ with $u=u_1,v=v_1$.
    Now let $u'_1=\lambda^{-1}\Psi^\top u$ and $v'_1=\pm\lambda^{-1}\Psi^\top v$ where the sign is so that $\omega_0(u'_1,v'_1)=1$.
    Extend this to a symplectic basis $u'_1,v'_1,\dots,u'_n,v'_n$ of $\bR^{2n}$.

    Define $\Phi,\Phi'\in\Sp(2n)$ to be the symplectic matrices which map the standard basis $e_1,\dots,f_n$ to $u_1,\dots,v_n$ and $u'_1,\dots,v'_n$ respectively.
    Then $A=(\Phi')^{-1}\Psi^\top\Phi$ satisfies
    $$ Ae_1 = \lambda e_1,\qquad Af_1 = \pm\lambda f_1 . $$
    Thus $A^\top$ maps the unit ball $B^{2n}(1)$ in $Z^{2n}(\lambda)$: if $\norm z\leq1$ then
    $$ \iprod{A^\top z,e_1}^2 + \iprod{A^\top z,f_1}^2 = \iprod{z,Ae_1}^2 + \iprod{z,Af_1}^2 = \lambda^2(\iprod{z,e_1}^2 + \iprod{z,f_1}^2) \leq \lambda^2 . $$
    Because $\lambda<1$, this means that $A^\top$ does not have the linear squeezing property.
    Because $\Phi,\Phi'$ are symplectic this means that $\Psi$ does not have the linear squeezing property.
    \qed

\eproof

Let $A\subseteq\bR^{2n}$ be an arbitrary subset.
Define its {\emphcolor linear symplectic width} to be
$$ w_L(A) = \sup\set{\pi r^2}[\psi(B^{2n}(r))\subseteq A\hbox{ for some $\psi\in\ASp(\bR^{2n})$}] . $$
The linear symplectic width has the following properties:
\benum
    \item {\emphcolor monotonicity}: if $\psi(A)\subseteq B$ for $\psi\in\ASp(\bR^{2n})$ then $w_L(A)\leq w_L(B)$.
    \item {\emphcolor conformality}: $w_L(\lambda A)=\lambda^2w_L(A)$.
    \item {\emphcolor nontriviality}: $w_L(B^{2n}(r))=w_L(Z^{2n}(r))=\pi r^2$.
\eenum
(1) is obvious, (2) is because $\psi(B^{2n}(r))\subseteq A$ iff $\psi(B^{2n}(\lambda r))\subseteq\lambda A$.
(3) is due to \refmath{affnonsqueeze}: $B^{2n}(r')$ embeds into $B^{2n}(r),Z^{2n}(r)$ iff $r'\leq r$ so $w_L(B^{2n}(r))=\pi r^2$.

\blemm

    Anti-symplectic linear maps have determinant $(-1)^n$.
    Furthermore anti-symplectic affine maps preserve linear symplectic width.

\elemm

\bproof

    If $\Psi^*\omega_0=-\omega_0$ then $\Psi^*\omega_0^n=(-1)^n\omega_0$.
    Because $\omega_0^n$ is a volume form, this means that $\det\Psi=(-1)^n$ as desired.

    For the second point note that the composition of two anti-symplectic maps is symplectic, of a symplectic with an anti-symplectic is anti-symplectic.
    Define $D=\diag(I,-I)$; this is anti-symplectic and $D^2=I$.
    Thus $\Psi$ is anti-symplectic iff $\Psi D$ is symplectic.
    Furthermore $D$ is orthonormal and thus $D(B^{2n}(r))=B^{2n}(r)$.

    This shows that $w_L(A)$ can be defined by taking the supremum over $\psi(B^{2n}(r))\subseteq A$ where $\psi$ is an anti-symplectic affine map.
    Thus $w_L$ is monotonic with respect to anti-symplectic affine maps, and we have shown the second claim.
    \qed

\eproof

An {\emphcolor ellipsoid} is a subset $E\subseteq\bR^{2n}$ of the form $\set{z\in\bR^{2n}}[Q(z)\leq1]$ where $Q$ is any positive definite quadratic form.
That is $Q(z)=z^\top\bar Qz$ for a symmetric positive-definite $\bar Q\in\Sym_+(2n)$.
Notice that if $\Psi\in\Mat_{2n}$ is nonsingular then $\Psi E$ is an ellipsoid as well:
$$ \Psi(E) = \set{z\in\bR^{2n}}[Q(\Psi^{-1}z)\leq1] , $$
and $Q\circ\Psi^{-1}$ is a positive definite quadratic form.

\bthrm

    Let $\Psi\in\Mat_{2n}$ be a linear map.
    Then the following are equivalent.
    \benum
        \item $\Psi$ preserves the linear symplectic width of all ellipsoids.
        \item $\Psi$ is either symplectic or anti-symplectic.
    \eenum

\ethrm

\bproof

    As shown above, symplectic and anti-symplectic maps preserve all widths, so (2) implies (1).
    So assume (2) and we will show that $\Psi$ has the linear nonsqueezing property.
    Let $B$ be a linear symplectic ball of radius $r$ and $Z$ a linear symplectic cylinder of radius $R$ such that $\Psi B\subseteq Z$.
    Then by hypothesis
    $$ \pi r^2 = w_L(B) = w_L(\Psi B) \leq w_L(Z) = \pi R^2 \implies r\leq R $$
    as desired.

    A symplectic ball cannot be embedded into a subspace of dimension $<2n$, so if $\Psi$ were singular then $w_L(\Psi B^{2n}(1))=0$, contradicting (1).
    Thus $\Psi$ is nonsingular, and $\Psi^{-1}$ satisfies (1) because 
    $$ w_L(\Psi^{-1}E) = w_L(\Psi\Psi^{-1}E) = w_L(E) $$
    as $\Psi^{-1}E$ is an ellipsoid.

    Thus both $\Psi$ and $\Psi^{-1}$ have the linear nonsqueezing property, so by \refmath{affrigid} $\Psi$ is symplectic or anti-symplectic.
    \qed

\eproof

\blemm

    Let $g\colon V\times V\to\bR$ be an inner product on a symplectic space.
    Then there exists a basis $u_1,\dots,u_n,v_1,\dots,v_n$ of $V$ which is both $g$-orthogonal and $\omega$-standard.
    More so, the basis can be chosen so that $g(u_j,u_j) = g(v_j,v_j)$ for all $j$.

\elemm

\bproof

    Consider $V=\bR^{2n}$ with the Euclidean inner product and assume that
    $$ \omega(z,w) = \iprod{z,Aw} . $$
    Then $A$ is nondegenerate and antisymmetric: $A^\top=-A$.
    Thus $iA\in\Mat_{2n}(\bC)$ is Hermitian, which means its eigenvalues are all real, so $A$'s eigenvalues are all pure imaginary.
    Furthermore the eigenvalues come in pairs $\pm i\alpha_j$ for $j=1,\dots,n$, with $\alpha_j>0$.

    As $iA$ is Hermitian, we can find $z_j=u_j+iv_j\in\bC^{2n}$ an orthonormal basis of eigenvectors:
    $$ Az_j = i\alpha_jz_j,\qquad \bar{z_j}^\top z_k = \delta_{jk} . $$
    And $\bar z_j$ fill out the basis of eigenvectors, so we have that $z_j^\top z_k=0$.
    Thus
    $$ Au_j = -\alpha_jv_j,\qquad Av_j = \alpha_ju_j,\qquad u_j^\top v_k = u_j^\top u_k = v_j^\top v_k = 0\hbox{ for $j\neq k$}. $$
    And so
    $$ \omega(u_j,v_j) = u_j^\top Av_j = \alpha_j\norm{u_j}^2 > 0, $$
    and similarly $\omega(u_j,v_k)=\omega(u_j,u_k)=\omega(v_j,v_k)=0$ for $j\neq k$.
    We can thus rescale $u_j,v_k$ to get the desired basis.

\eproof

Let $r=(r_1,\dots,r_n)$ be an $n$-tuple with $0<r_1\leq\cdots\leq r_n$, and define the ellipsoid
$$ E(r) = \set{z\in\bC^n}[\sum_{i=1}^n\abs{\frac{z_i}{r_i}}^2\leq1] . $$

\blemm

    Any ellipsoid
    $$ E = \set{w\in\bR^{2n}}[w^\top Aw\leq1] $$
    is symplectomorphic to a unique $E(r)$.

\elemm

\bproof

    We can write $E$ as $\set{w\in\bR^{2n}}[g(w,w)\leq1]$ for the inner product
    $$ g(u,v) = u^\top Av . $$
    By the above lemma there is a basis $u_1,\dots,u_n,v_1,\dots,v_n$ of $\bR^{2n}$ which is both symplectic and orthogonal for $g$.
    Define $r_j$ so that $g(u_j,u_j)=g(v_j,v_j)=1/r_j^2$.
    We may rearrange the $u_j,v_j$ so that the $r_j$ are ascending; this does not affect them forming a symplectic basis.

    Let $\Phi\colon\bR^{2n}\to\bR^{2n}$ be the symplectomorphism mapping the standard basis to this new basis:
    $$ \Phi(z) = \sum_{j=1}^n(x_ju_j + y_jv_j),\qquad z=(x_1,\dots,x_n,y_1,\dots,y_n) . $$
    Thus
    $$ g(\Psi z,\Psi z) = \sum_{j=1}^n\frac{x_j^2+y_j^2}{r_j^2} , $$
    which means that
    $$ \Psi^{-1}E = E(r) , $$
    as desired.

    Now to show uniqueness of $r=(r_1,\dots,r_n)$, define the matrix
    $$ \Delta(r) = \diag(1/r_1^2,\dots,1/r_n^2,1/r_1^2,\dots,1/r_n^2) $$
    so that
    $$ E(r) = \set{z\in\bR^{2n}}[z^\top\Delta(r)z\leq1] . $$
    Because $\Psi E(r)$ maps to the ellipsoid defined by $\Psi^\top\Delta(r)\Psi$, we want to show that if $\Psi^\top\Delta(r)\Psi=\Delta(r')$, then $r=r'$.
    We know that $J_0\Psi^\top=\Psi^{-1}J_0$ so
    $$ \Psi^{-1}J_0\Delta(r)\Psi = J_0\Delta(r') . $$
    And so $J_0\Delta(r)$ and $J_0\Delta(r')$ have the same eigenvalues.
    But these matrices correspond to the complex matrices $i\diag(1/r_1^2,\dots,1/r_n^2)$ and $i\diag(1/{r'_1}^2,\dots,1/{r'_n}^2)$.
    Thus we must have $r=r'$ as desired.
    \qed

\eproof

Note that $E(r,\dots,r)$ is just $B^{2n}(r)$.

\bnote

    In the case of $n=1$, symplectic forms are volume forms and symplectomorphisms are the same as area-preserving maps.
    Thus the above lemma states that every ellipse in $\bR^2$ can be mapped into a circle by an area-preserving linear transformation, and that this circle is unique.

\enote

Define the {\emphcolor symplectic spectrum} of an ellipsoid $E$ to be the unique $n$-tuple $r=(r_1,\dots,r_n)$ with $0<r_1\leq\cdots\leq r_n$ such that $E$ is symplectomorphic to
$E(r)$.
This is an invariant which uniquely determines the ellipsoid up to symplectomorphism.
So if $E$ has symplectic spectrum $r$, we see that
$$ \vol E = \vol E(r) = \frac{\pi^n}{n!}\prod_{j=1}^nr_j^2 $$
where the last equality is just an explicit integral.

\bthrm

    Let $E\subseteq\bR^{2n}$ be an ellipsoid, then
    $$ w_L(E) = \sup_{B\subseteq E}w_L(B) = \inf_{Z\supseteq E}w_L(E) , $$
    where the supremum runs over all linear (or affine) symplectic balls contained in $E$, and the infimum over all linear symplectic cylinders containing $E$.
    Let $r=(r_1,\dots,r_n)$ be $E$'s symplectic spectrum, then
    $$ w_L(E) = \pi r_1^2 . $$

\ethrm

\bproof

    Suppose $\Psi\in\Sp(2n)$ is a symplectomorphism such that $\Psi E=E(r)$.
    Then we know
    $$ B^{2n}(r_1) \subseteq E(r) \subseteq Z^{2n}(r_1) \implies \Psi^{-1}B^{2n}(r_1) \subseteq E \subseteq \Psi^{-1}Z^{2n}(r_1) . $$
    The widths of the bounds are both $\pi r_1^2$ so we have
    $$ w_L(E) = \pi r_1^2 . $$
    And furthermore we then get
    $$ \inf_{Z\supseteq E}w_L(Z) \leq \pi r_1^2 \leq \sup_{B\subseteq E}w_L(B) $$
    because $\Psi^{-1}Z^{2n}(r_1)$ is in the infimum on the left, and $\Psi^{-1}B^{2n}(r_1)$ is in the supremum on the right.
    But if $B\subseteq E$ is a symplectic ball and $E\subseteq Z$ is a symplectic cylinder, then $w_L(B)\leq w_L(Z)$ by monotonicity, and so every element in the supremum
    is less than every element in the infimum, so we have equality.
    \qed

\eproof

For an ellipsoid $E$, define its {\emphcolor dual} to be
$$ E^* = \set{v\in\bR^{2n}}[\iprod{v,e}\leq1\hbox{ for all $e\in E$}] . $$
Let us temporarily denote the ellipsoid determined by the positive definite quadratic form $A$ by $E(A)$, and the inner product $g_A(v,u)=v^\top Au$.
Then we claim $E(A)^*=E(A^{-1})$.
Notice that $v\in E(A^{-1})$ iff $v^\top A^{-1}v\leq1$ iff $A^{-1}v\in E(A)$.
Thus for all $e\in E(A)$ we have by Cauchy-Schwarz
$$ \iprod{v,e} = g_A(A^{-1}v,e) \leq 1 $$
so $v\in E(A)^*$.
Conversely if $v\in E(A)^*$ then
$$ \iprod{v,e} = g_A(A^{-1}v,e) \leq 1 $$
for all $e\in E(A)$.
Let $\alpha=\sqrt{g_A(A^{-1}v,A^{-1}v)}$ be $A^{-1}v$'s norm, so that $A^{-1}v/\alpha$ has $g_A$-norm $1$ and thus is in $E(A)$.
But
$$ g_A(A^{-1}v,A^{-1}v/\alpha) = \alpha \leq 1 $$
meaning $A^{-1}v\in E(A)$ and thus $v\in E(A^{-1})$ as desired.

Note that $A$ being positive definite means $A^{-1}$ is too, so $E(A^{-1})=E(A)^*$ is an ellipsoid.
Furthermore it is immediate from this that $E^{**}=E$ and $(\Psi E)^*=(\Psi^\top)^{-1}E^*$, the latter being because
$$ (\Psi E(A))^* = E((\Psi^\top)^{-1}A\Psi^{-1})^* = E(\Psi A^{-1}\Psi^\top) = (\Psi^\top)^{-1}E(A)^* . $$

